%%=============================================================================
%% Conclusie
%%=============================================================================

\chapter{Conclusie}%
\label{ch:conclusie}

% TODO: Trek een duidelijke conclusie, in de vorm van een antwoord op de
% onderzoeksvra(a)g(en). Wat was jouw bijdrage aan het onderzoeksdomein en
% hoe biedt dit meerwaarde aan het vakgebied/doelgroep? 
% Reflecteer kritisch over het resultaat. In Engelse teksten wordt deze sectie
% ``Discussion'' genoemd. Had je deze uitkomst verwacht? Zijn er zaken die nog
% niet duidelijk zijn?
% Heeft het onderzoek geleid tot nieuwe vragen die uitnodigen tot verder 
%onderzoek?

Dit laatste hoofdstuk van de batchelorproef bespreekt de bijdragen  
aan het onderzoeksdomein, reflecteert kritisch op de resultaten en 
formuleert een antwoord op de centrale onderzoeksvraag.
Aansluitend worden openstaande vragen en aanbevelingen voor verder 
onderzoek geformuleerd.

\section{Antwoord op de onderzoeksvraag}%
\label{sec:conclusie-antwoord}

Vanuit de centrale onderzoeksvraag van deze bachelorproef kunnen enkele belangrijke conclusies worden getrokken. 
Deze onderzoeksvraag werd als volgt geformuleerd:
Hoe kan een competentieframework, in combinatie met een bijhorend leerplatform,
bijdragen aan een effectievere instroom en onboarding van startende PL/I en CICS ontwikkelaars binnen een z/OS omgeving?

Op basis van de literatuurstudie en de interviews met professionals 
bij Euroclear, in samenhang met de ontwikkeling en validatie van zowel het framework 
als de proof of concept kan hierop een positief antwoord worden 
geformuleerd.

Een rolspecifiek en gevalideerd competentieframework biedt een structuur die de groei van opleiding naar het werkveld verkort.
Het maakt het aanleren van complexe vakkennis toegankelijker voor starters en geeft hun een duidelijk groeipad.
Voor organisaties biedt het een gestructureerde basis voor hun onboarding en interne opleidingstrajecten.

Een framework in samenhang met een leerplatform zorgt voor een nog grotere meerwaarde aan hun leer traject.
Studenten kunnen direct hun theoretishe kennis toepassen op gerichte oefeningen en kunnen zichzelf evalueren op hun vooruitgang.
Dit zorgt voor een veel gestuctureerdere manier van leren waar weinig tot geen buitenstaande interventie nodig is.

\section{Bijdrage aan het onderzoeksdomein}%
\label{sec:conclusie-bijdrage}

De bijdrage van de batchelorproef aan het onderzoeksdomein kan gesplitst worden in twee delen.

Theoretisch gericht werd er een duidelijk overzicht opgebouwd met de relevante literatuur rond het skillsgap in de mainframe industrie en de onboarding van ontwikkelaars in legacy omgevingen.
Dit overzicht toonde aan dat er al bestaande frameworks zijn rond legacy systemen maar dat deze niet gericht zijn op de combinatie PL1 en CICS development op z/OS systemen ontwikkelaarsrol.
Deze batchelorproef vult deze leegte dan ook aan met een gevalideerd rolspecifiek voor PL1 en CICS ontwikkelaars framework dat bestaat uit dertien kern competenties.

Praktish gericht werd er een proof of concept ontwikkeld die aantoont dat het competentie framework in de praktijk kan toegepast worden als doelstellingen van een educatieve website.
Het leerplatform kan ook direct gebruikt worden door Hogent studenten voor het aanleren van de belangerijkste compententies, maar ook als onboarding tool voor organisaties zoals Euroclear.

\section{Kritische reflectie}%
\label{sec:conclusie-reflectie}

Voor een compleet beeld te krijgen van de resultaten is een kritische reflectie een heel belangerijk aspect.
De initele resultaten van de bachelorproef zijn positief maar enkele beperkingen kunnen een verkeerd beeld geven.

\subsection{beperkingen onderzoek}%
\label{sec:beperkingen-reflectie}

Een eerste beperking die vast kan gesteld worden is de scope van de interviews bij professionals.
De geinterviewde professionals zijn allemaal van één zelfde organisatie genaam Euroclear, hierdoor is het framework erg sterk gebaseerd op de verwachtingen van deze organisatie.
Andere organisaties in verschilende sectoren kunnen andere verwachtingen stellen voor startende ontwikkelaars op hun systemen.
Een mogelijke verbeteringen is het stellen van deze vragen bij vershillende organisatie zodat er een meer globaal beeld kan gevormd worden.

Een tweede bepreking is dat de effectiviteit van het leerplatform slecht beperkt kan worden getoets.
De leersite word gevalideerd door professionals en slechts enkele studenten over een gelimiteerde tijdsperiode.
Hierdoor kan er geen zeer concrete conclusie getrokken worden uit deze testen.
Dit zorgt er voor dat de vraag open blijft of het leerplatform de studenten effectief en zelfstandig door de leerstof begeleid enof externe ondersteuning nodig is.
Deze vraag zou kunnen opgelost worden door het valideren van de educatieve website bij student van Hogent die de opleiding Mainframe Expert volgen.

\subsection{verwacht resultaat onderzoek}%
\label{sec:verwachte-resultaat-reflectie}

Het verwachte resultaat van deze batchelorproef lopen in lijn met de resultaten van het onderzoek.
Er werd bevestigd dat de nood aan een rolspecifiek framework en een zeker skillsgap zeker aanwezig zijn, wat niet heel verrassend is want bestaande literatuur wijst op deze duidelijke signalen.
Wat meer opviel tijdens de intervieuws was de nadruk die professionals legden op de historishe en conceptuele kennis.
De softskills naast de hardskills werden als zekere succesfactoren aangekaart voor nieuwe starters.
Dit aspect werd in bestaande framework te wienig besproken en is een duidelijk werkpunt voor toekomstige frameworks.

\section{Aanbevelingen verder onderzoek}%
\label{sec:conclusie-aanbevelingen}


Deze batchelorproef zet ook aan tot opvolgend onderzoek door enkele nieuwe vragen die naar boven zijn gekomen.

Een zeer belangerijke vraag die kan gesteld worden is welke impact het framework en leerplatform hebben op de effectieviteit van onboarding binnen legacy development teams.
Door studenten van Hogent gedurende een volledig acadamiejaar te laten experimentern met het leerplatform en vervolgens hun onboarding process op te volgen binnen hun organisatie.
Met deze gegevens kan gericht nagegaan worden of het platform de overgang naar een mainframe omgeving heeft verkort en de productiviteit van starters heeft verhoogt.

Naast deze hoofdvraag kan er ook nog een tweede vraag gesteld worden, deze is gebaseerd op de vraag of externe begeleiding nodig is.
Het integreren van een AI assistent binnen het leerplatform zou deze rol van externe begeleider kunnen vullen.
Met de recente vooruitgang in taalmodellen rond mainframe applicaties zou een AI assitent in staat zijn om intelligente feedbackt te geven op praktishe oefeningen, en de student begeleiden bij het debuggen van hun code zoals een ervaren begeleider dat zou doen.
Dit zou de zelfstandigheid van de student kunnen vergroten en de drempel tot fouten maken ook verkleinen.


Tot slot kan geconcludeerd worden dat het ontwikkelen van gestructureerde frameworks en opleidingstrajecten een noodzaak zijn voor het beperken van het vergrijzings probleem binnen z/OS ontwikkeling.
Organisaties die vandaag investeren in het implementeren van deze frameworks bouwen een kennis basis op bij nieuwe starters die morgen de kritieke infrastructuur draaiende houden.







